\chapter{Conclusions}
This chapter encapsulates the outcomes and insights gathered from the implementation of the ClassiFish system—an end-to-end fish species detection and feedback-enabled mobile application powered by YOLOv11. It revisits the original aims, presents the study's conclusions, outlines its contributions, critiques the process, and suggests potential directions for future work.

\section{Revisiting the Aims and Objectives}
The primary aim of this study was to develop a fish identification system using the YOLO-v11 model. The system was intended to identify various fish species from images and allow users to submit new fish images and provide feedback on identification results. To meet this aim, several specific objectives were pursued. First, images of different fish species were collected from local wet markets. These images were manually annotated to build a labeled dataset. Using this dataset, a YOLO-v11 model was developed and trained to recognize and classify 18 different fish species. The performance of the model was evaluated using metrics such as mean Average Precision (mAP). A mobile application was also created to make the identification system accessible to users. This app not only displays the predicted fish species but also provides related information such as the species’ name, taxonomic classification, physical description, nutritional content, conservation status, and health benefits. Additionally, the app includes a feedback feature that enables users to confirm or correct the species identified by the system.

\section{Conclusions}
ClassiFish successfully demonstrated that deep learning-based object detection (YOLOv11) can be effectively deployed on mobile devices for real-time fish classification. With a dataset of 2,700 raw and over 9,000 augmented images, the model achieved a precision of 0.72, recall of 0.68, F1-score of 0.70, and mAP@0.5 of 0.75—indicating acceptable performance, particularly for clearly distinguishable fish classes. The app includes a feedback mechanism that allows users to submit corrections, which can be reviewed for potential use in retraining the model.

\section{Contributions}
This study presents several key contributions aimed at advancing ecological education and fisheries research. First, it provides a localized and annotated dataset featuring 18 Philippine fish species, ensuring accurate recognition and classification for related applications. Additionally, a YOLOv11 model has been trained and optimized for mobile inference by converting it to TensorFlow Lite, enabling efficient deployment on mobile devices. A fully functioning Flutter-based mobile application has been developed, integrating real-time detection, informative species display, and user feedback submission to enhance interaction and usability. Furthermore, a Django-powered web admin panel facilitates the management and validation of user-submitted corrections, ensuring continuous model refinement through future retraining. Lastly, the study establishes a community-driven model improvement loop that fosters collaboration, supports ecological education, and contributes to fisheries research by refining detection accuracy based on user input and expert validation. These contributions collectively enhance the accessibility and reliability of AI-powered fish species identification in a practical and user-centric manner.

\section{Critique and Limitations}
Despite its achievements, ClassiFish has several limitations:
\begin{itemize}
    \item The model struggled with visually similar fish species such as Tamban and Pulag-ikog, likely due to inter-class similarity and limited variance in image conditions.
    \item The image dataset, while diverse, was limited to natural lighting conditions and wet market scenes.
    \item The feedback mechanism relies on user input which may be inconsistent or inaccurate if not reviewed properly.
\end{itemize}

\section{Future Work}
To further enhance this study, several recommendations are proposed. First, expanding the dataset to include a broader range of fish species and varying lighting conditions would improve generalization, ensuring robust performance across diverse real-world scenarios. Additionally, implementing an active learning framework would allow verified user feedback to directly influence retraining, reducing the need for full manual intervention and accelerating model improvement. Finally, incorporating gamification elements within the mobile application would encourage user engagement, fostering increased participation and interaction while promoting ecological awareness and fisheries research. These future directions aim to refine and extend the study’s impact, making AI-powered fish species identification more effective and user-centric.

\section{Final Remarks}
ClassiFish bridges the gap between technology and marine biodiversity awareness through an accessible, intelligent system. This study not only contributes to computer vision applications in aquaculture and fisheries but also showcases the impact of community involvement in artificial intelligence systems. The road ahead offers exciting possibilities to further empower users and researchers alike in sustainable marine practices.